\documentclass[10pt,a4paper]{amsart}
\usepackage[margin=19mm]{geometry}
\usepackage{amsmath,amssymb,mathtools}
\usepackage[scr=rsfs,cal=euler]{mathalfa}
\usepackage{xfrac}
\usepackage[unicode,hidelinks,bookmarks=false]{hyperref}
\newcommand{\ie}{i.e.\ }
\newcommand{\cf}{cf.\ }
\newcommand{\resp}{respectively\ }
\newcommand{\Subsection}{\subsection}
\providecommand{\nobreakdash}{\nobreak\hbox{-}}
\setlength{\emergencystretch}{2em}
\multlinegap0pt
\usepackage{bm}
\usepackage{bbm}
\usepackage{dsfont}
\usepackage{xspace}
\usepackage{cancel}
\usepackage{mathtools}
\usepackage{amsthm}
\makeatletter
\@ifundefined{theorem}{\newtheorem{theorem}{Theorem}}{}
\@ifundefined{assumption}{\newtheorem{assumption}[theorem]{Assumption}}{}
\@ifundefined{lemma}{\newtheorem{lemma}[theorem]{Lemma}}{}
\@ifundefined{proposition}{\newtheorem{proposition}[theorem]{Proposition}}{}
\makeatother
\newcommand{\E}{\mathbb{E}}
\newcommand{\Ohp}{O_{v.h.\p}}
\newcommand{\Op}{O_{\p}}
\newcommand{\centered}{C_n}
\newcommand{\eps}{\varepsilon}
\newcommand{\ohp}{o_{v.h.\p}}
\newcommand{\one}{\textup{\textbf{1}}}
\newcommand{\p}{\mathbb{P}}
\newcommand{\rhs}{{r.h.s.~}}
\newcommand{\scaling}{s_n}
\newcommand{\wconv}{\xrightarrow[\textit{n} \to +\infty]{\quad {w} \quad}}
\newcommand{\w}{\boldsymbol{w}}
\newcommand{\x}{{{v}}^+}
\newcommand{\y}{{{v}}^-}
\title[Spectrum of Directed Inhomogeneous Random Graphs]{Spectrum of Directed Inhomogeneous Random Graphs}
\author{Rajat Subhra Hazra, Giacomo Passuello}
\date{Source: arXiv:2607.08696v1 (CC BY 4.0)}
\begin{document}
\maketitle

\noindent\emph{Source and licence.} Spectrum of Directed Inhomogeneous Random Graphs. Version arXiv:2607.08696v1 (CC BY 4.0), distributed under the Creative Commons Attribution 4.0 International licence, as recorded in the arXiv licence metadata for this exact version and shown on the abstract page https://arxiv.org/abs/2607.08696v1. Full rights notice: licenses/arxiv-2607.08696-CC-BY-4.0.txt.\par\smallskip

\noindent\textit{Editorial scope.} Counted unit: the Lemma whose source label is \texttt{lem:fixed-point} in the article (source0.tex lines 1700--1705, source label \texttt{lem:fixed-point}), delivered together with the article's own complete original proof of it (source0.tex lines 1707--1742, one \texttt{proof} environment of 36 lines). No mathematics is written, repaired or re-derived: statement and proof are the article's own text line for line. Required context retained uncounted: eq:sp-chung-lu (lines 277--281); ass:rank-one (lines 285--300); assumption:xi (lines 437--439); thm:LLN-A (lines 440--446); lem:5 (lines 1329--1331); lem:1 (lines 1345--1375); eq:k=1 (lines 1351--1352); prop:prop-W (lines 1411--1417); eq:lambda (lines 1684--1687). Every definition, display and label of those lines is retained unchanged. Omitted from this package and not redistributed: 1--276, 282--284, 301--436, 447--1328, 1332--1344, 1353--1410, 1418--1683, 1688--1699, 1706--1706, 1743--1995. Nothing inside a retained excerpt is omitted or altered: every retained body is delivered exactly as the article's lines are written, so no omission receipt of any kind accompanies this package. Citations: the retained excerpts carry no citation command at all, so no bibliography entry of the article is transcribed and no reference is redistributed. Reference adaptation: none was needed and none was made; every reference inside the delivered unit resolves inside it, so no unresolved reference or citation is produced. Printed slips kept literal: no printed word, symbol, hypothesis or number of the counted statement or of its proof was altered. Layout and class: the article's own class is \texttt{amsart} with packages including babel, enumitem, csquotes, bbm, subcaption, float, fancyhdr, etoolbox, xcolor, hyperref, amsmath, amssymb, mathrsfs, palatino; the wrapper is the lane's pinned \texttt{amsart} editorial wrapper at 10pt on A4, which restates the theorem-like environments the retained text needs and adds the article's own macro definitions that the retained text needs (\texttt{\textbackslash E}, \texttt{\textbackslash Ohp}, \texttt{\textbackslash Op}, \texttt{\textbackslash centered}, \texttt{\textbackslash eps}, \texttt{\textbackslash ohp}, \texttt{\textbackslash one}, \texttt{\textbackslash p}, \texttt{\textbackslash rhs}, \texttt{\textbackslash scaling}, \texttt{\textbackslash w}, \texttt{\textbackslash wconv}, \texttt{\textbackslash x}, \texttt{\textbackslash y}), with extra wrapper packages (bm, bbm, dsfont, xspace, cancel, mathtools). The delivered numbering is the unit's own. Assets: no retained span contains a figure environment, a graphics-inclusion command, a PDF-page inclusion, a table or a TikZ picture; no image file is copied into this package, the package declares no asset dependency and no image file is redistributed with it. Licence: version arXiv:2607.08696v1 of this article is distributed under the Creative Commons Attribution 4.0 International licence, as recorded in the arXiv licence metadata for this exact version and shown on the abstract page https://arxiv.org/abs/2607.08696v1. A full rights notice is delivered at licenses/arxiv-2607.08696-CC-BY-4.0.txt. Characters: every retained excerpt is pure ASCII, so the delivered engine, XeLaTeX with its default Latin Modern text font, reports no missing character.
\begin{equation}
\label{eq:sp-chung-lu}
        p_{xy}:=\mathbb P(a_{xy}=1)=\frac{s_n w_x^+w_y^-}{\w},
        \qquad x,y\in[n].
\end{equation} 

\begin{assumption}[Rank-one weight profile]
\label{ass:rank-one}
There exist constants $0<c<C<\infty$ such that
\begin{equation}
\label{eq:weight-bounds}
        c\le w_x^\pm\le C,\qquad x\in[n],\ n\ge1.
\end{equation}
Moreover there exists a compactly supported probability measure
$\rho=\rho_{w^+,w^-}$ on $(0,\infty)^2$ such that
\begin{equation}
\label{eq:rho-one}
        \frac1n\sum_{x=1}^n
        \delta_{\sqrt{n/\w_n}(w_x^+,w_x^-)}
        \wconv \rho .
\end{equation}
\end{assumption}

\begin{assumption}\label{assumption:xi} 
There exists $\xi>4$ such that $ \log^\xi(n)\ll \scaling \ll n$.
\end{assumption}

\begin{theorem}[Existence of outlier - rank-one case] \label{thm:LLN-A}
Consider the Chung--Lu model with $p_{xy}$ as in Eq.\ \eqref{eq:sp-chung-lu}.
If Assumptions \ref{ass:rank-one} and \ref{assumption:xi} hold, then there exists $K>0$ and $\eta>1$ such that
\begin{equation}
    \p\left(\max\left\{\left|\lambda_1(A_n)- \lambda_1(\E[A_n]) \right|,\max_{ 2 \le j\le n}|\lambda_j(A_n)|\right\}\ge K\sqrt{\scaling} \right)\le e^{-(\log(n))^{\eta}}.
\end{equation}
\end{theorem}

\begin{lemma}  \label{lem:5}
        $| (\y)^t \centered \x | = \Op\left( \sqrt{\frac{\scaling}{n}}\right).$
\end{lemma}

\begin{lemma} \label{lem:1}
There exists a constant $K_1<+\infty$ such that, for $2\le k \le L$,
\begin{equation}
    \left|\E \left[(\y)^t\centered^k\x\right]\right| \le \left(K_1 \scaling\right)^{k/2}.
\end{equation}
Moreover, 
\begin{equation} \label{eq:k=1}
(\y)^t \centered \x = \ohp( \scaling).\end{equation}
\begin{proof}
By Proposition \ref{prop:prop-W}, there exists $K>0$ and $\eta>1$ such that for the event $\mathcal{A}\coloneq \{\|\centered\|\le C\sqrt{\scaling}\}$ it holds $\p(\mathcal{A}^\complement)\le e^{-(\log(n))^\eta}$. Then
\begin{equation}
    \left|\E\left[(\y)^t \centered^k \x \right]\right|\le     
    \left|\E\left[(\y)^t \centered^k \x \one_{\mathcal{A}} \right]\right|
    + \left|\E\left[(\y)^t \centered^k \x \one_{\mathcal{A}^\complement} \right]\right|.
\end{equation}
The first term is bounded by  $\|v^-\|\|v^+\|\E[\|\centered\|^k\one_{\mathcal A}] \le K_1\scaling^{{k}/2}$. For the second term, note that $\left((\y)^t \centered^k \x\right)^2$ is bounded by a power of $n$, say $n^{kC'}$, for $C'>0$. Hence by Cauchy--Schwartz,
\begin{align*}
    \left|\E\left[(\y)^t \centered^k \x \one_{\mathcal{A}^\complement} \right]\right|
    &  \le \left(\E\left[\left((\y)^t \centered^k \x\right)^2\right]\right)^{\frac12}\p\left({\mathcal{A}^\complement}\right)^{\frac12} \\
    & \le n^{kC'/2}e^{-(\log(n))^\eta/2}=o(1).
\end{align*}
To prove \eqref{eq:k=1}, recall that $\centered=A_n-\E[A_n]$. Hoeffding  inequality gives,  for every $\eps>0$,
\begin{equation}
\begin{split}
    \p
    \Bigg(\Bigg|\sum_{x,y \in [n]}v_x^- a_{xy} v_y^+-\sum_{x,y \in [n]}&\E\left[v_x^- a_{xy} v_y^+\right]\Bigg|> \eps\scaling\Bigg)\\ &\le 2 \exp \left(-\frac{2\eps^2 \scaling^2}{n^2 ((\max_{x} w_x^\pm-\min_x w_x^\pm)/\w)^2}\right).
    \end{split}
\end{equation}
Since weights are bounded, the \rhs is at most $O(\exp(-2\eps^2(\log(n))^{2\xi}))$.
\end{proof}
\end{lemma}

\begin{proposition}
\label{prop:prop-W}
If Assumptions \ref{ass:rank-one} and \ref{assumption:xi} hold, there exist $K_1>0$ and $\eta>1$ such that, for large $n$,
\begin{equation}
    \p\left(\|\centered\|\ge K_1\sqrt{\scaling}\right)\le e^{-(\log(n))^{\eta}}.
\end{equation}
\end{proposition}

\begin{equation} \label{eq:lambda}
    \lambda 
    = (\y)^t\x \scaling + \scaling \sum_{k=2}^{L} (v^-)^t \left(\frac{\E[\centered]}{\lambda}\right)^k v^+ + \Ohp\left(\sqrt{\frac{\log(n)^{\xi}}{n}}\right)
\end{equation}

\begin{lemma} \label{lem:fixed-point}
   In the present setting, it holds
     % \begin{equation} \label{eq:lambda-3}
        $\lambda -\tilde{\lambda}=\scaling\frac{(v^-)^t \centered v^+}{\tilde\lambda}+\ohp\left(\sqrt{\frac{\scaling}{n}}\right).$
    % \end{equation}
\end{lemma}

\begin{proof}
Combining \eqref{eq:lambda} with the definition of $\tilde \lambda$ we get
\begin{equation} \label{eq:equation-H^{(3)}_n}
    \lambda-\tilde\lambda
    = \scaling\frac{(v^-)^t \centered v^+}{\lambda}+\scaling\sum_{k=0}^{\log(n)}\left(\frac{1}{\phantom{\tilde\lambda} \!\!\!\!\ \lambda^k}-\frac{1}{\tilde\lambda^k}\right) (v^-)^t {\E[C^k_n]} v^+ + \Ohp\left(\sqrt{\frac{\log(n)^{\xi}}{n}}\right).
\end{equation}
By Theorem \ref{thm:LLN-A} it holds
\begin{equation} \label{eq:frac1}
    \frac{1}{\phantom{\tilde\lambda} \!\!\!\!\ \lambda^k}-\frac{1}{\tilde\lambda^k}= (\tilde\lambda-\lambda)\left(\frac{\sum_{j=1}^{k-1}\tilde\lambda^j \lambda^{k-j}}{\lambda^k\tilde\lambda^k}\right)=(\tilde\lambda-\lambda)\Ohp\left(\frac{k}{\scaling^{k+1}}\right),
\end{equation}
so that, applying Lemma \ref{lem:1}, it results
\begin{equation} \label{eq:frac2}
\begin{split}
    \left|\scaling\sum_{k=0}^{\log(n)}\left(\frac{1}{\phantom{\tilde\lambda} \!\!\!\!\ \lambda^k}-\frac{1}{\tilde\lambda^k}\right) (v^-)^t {\E[C^k_n]} v^+\right|
    & \le |\lambda-\tilde\lambda|\, \Ohp\left( \scaling \sum_{k=2}^{\log(n)} \frac{k}{\scaling^{k+1}} (K_1 \scaling)^{k/2}\right)\\
    & =\Ohp\left(\frac{|\lambda-\tilde\lambda|}{\scaling}\right),
\end{split}
\end{equation}
where we used that $\sum_{k=2}^{\log(n)} k/\scaling^{k/2}= O(1/\scaling)$.
As a consequence,
\begin{equation} \label{eq:sp-display}
    \lambda-\tilde\lambda=\scaling\frac{(v^-)^t \centered v^+}{\lambda} + \Ohp\left(\frac{|\lambda-\tilde\lambda|}{\scaling}\right)+ \Ohp\left(\sqrt{\frac{\log(n)^{\xi}}{n}}\right).
\end{equation}
By Lemma \ref{lem:5} and Theorem \ref{thm:LLN-A} it holds $\scaling\frac{(v^-)^t {\centered} v^+}{\lambda} = \Ohp(\sqrt{\frac{\scaling}{n}})$. Then, Eq.\ \eqref{eq:sp-display} implies 
\begin{equation} \label{eq:lambdas}
|\lambda-\tilde\lambda|=  \Ohp\left(\sqrt{\frac{\scaling}{n}}\right).\end{equation}
Consequently, we can omit the second addend in the r.h.s.\ of Eq.\ \eqref{eq:sp-display} and get the more precise estimate
    \begin{equation} \label{eq:lambda-2}
        \lambda -\tilde{\lambda}=\scaling\frac{(v^-)^t \centered v^+}{\lambda}+\ohp\left(\sqrt{\frac{\scaling}{n}}\right).
    \end{equation}
 Reasoning as in \eqref{eq:frac1}, and thanks to \eqref{eq:lambdas}, we also have 
 \begin{equation}
     \left|\left(\frac{1}{\phantom{\tilde\lambda} \!\!\!\!\!\ \lambda}-\frac{1}{\tilde\lambda}\right)(v^-)^t \centered v^+\right| = \ohp(|\lambda-\tilde\lambda|)=\ohp\left(\sqrt{\frac{\scaling}{n}}\right).
 \end{equation} 
Then, we can change the $\lambda$ in the denominator of \eqref{eq:lambda-2} to $\tilde \lambda$. This proves Lemma \ref{lem:fixed-point}.
\end{proof}

\end{document}
